%HOMEWORK template for M 325K
\documentclass{article}
\headheight=7pt         \topmargin=14pt
\textheight=574pt       \textwidth=445pt
\oddsidemargin=18pt     \evensidemargin=18pt

%Begin package import

\usepackage{amsmath, amssymb, amsthm}
\usepackage{mathtools}
\usepackage{pdfpages}
\usepackage{tikz-cd}

%End package import
%Begin custom commands

\newcommand{\C}{\mathbb{C}}
\newcommand{\R}{\mathbb{R}}
\newcommand{\Q}{\mathbb{Q}}
\newcommand{\Z}{\mathbb{Z}}
\newcommand{\N}{\mathbb{N}}
\newcommand{\F}{\mathbb{F}}
\newcommand{\abs}[1]{\left\lvert #1\right\rvert}
\newcommand{\RM}[1]{\mathrm{#1}}
\newcommand{\BF}[1]{\textbf{#1}}
\newcommand{\e}{\epsilon}
\newcommand{\ceil}[1]{\lceil{#1}\rceil}
\newcommand{\floor}[1]{\lfloor{#1}\rfloor}
\newcommand{\rarr}[0]{\rightarrow}
\newcommand{\IM}[1]{\text{Im}(#1)}
\newcommand{\Hom}[0]{\text{Hom}}
\newcommand{\Pow}[1]{\text{Pow}(#1)}
\newcommand{\angles}[1]{\langle #1 \rangle}

%End custom commands
%Begin custom theorems

\newtheorem{innercustomgeneric}{\customgenericname}
\providecommand{\customgenericname}{}
\newcommand{\newcustomtheorem}[2]{%
\newenvironment{#1}[1]
{%
\renewcommand\customgenericname{#2}%
\renewcommand\theinnercustomgeneric{##1}%
\innercustomgeneric%
}
{\endinnercustomgeneric}
}

\newcustomtheorem{THRM}{Theorem}
\newcustomtheorem{LEM}{Lemma}
\newcustomtheorem{EX}{Exercise}
\newcustomtheorem{COR}{Corollary}
\newcustomtheorem{PROB}{Problem}
\newcustomtheorem{claim}{Claim}

\newenvironment{solution}
{\renewcommand\qedsymbol{$\blacksquare$}\begin{proof}[Solution]}
{\end{proof}}

\newenvironment{definition}
{\renewcommand\qedsymbol{$\blacksquare$}\begin{proof}[Definition]}
{\end{proof}}

%End custom theorems
\begin{document}

\title{Chapter 2}
\author{Arham Lodha}%put your name here
\date{May 08, 2024}%enter the due date here
\maketitle

\section*{Topological Spaces}



% \begin{PROB}{3}\label{Problem 3.}
%     Show that the collection $\mathcal{T}_c$ given in Example 4 of §$12$ is a topology on the set $X$. Is the collection
%     $\mathcal{T}_\infty := \left\{ U \mid X - U \text{ is infinite or empty or all of } X \right\}$?    
% \end{PROB}
% \begin{proof}
%     Rule 1 is satisfied bc $X, \emptyset \in \mathcal{T}_c$. The union of a finite subcollection of countable subsets is countable or is X. The intersection of a finite subcollection of countable subsets is countable or X. Thus $\mathcal{T}_c$ is a topology. 
    
%     The intersection of a finite subcollection of $\mathcal{T}_\infty$, may be finite. For example Let $X = \Z_+$. $2\N, 3\N \cup \left\{ 2 \right\}\in \mathcal{T}_\infty $, but $2\N \cap \left(3\N \cup \left\{ 2 \right\}\right) = \left\{ 2 \right\} \not \in \mathcal{T}_\infty$. Thus there is a contradiction with rule 3.     
% \end{proof}
% \bigskip

% \begin{PROB}{4a}\label{Problem 4a.}
%     If $\left\{ \mathcal{T}_\alpha \right\}$ is a family of topologies on $X$ , show that $\bigcap \mathcal{T}_\alpha$ is a topology on $X$. Show that $T = \bigcup \mathcal{T}_\alpha$ is a topology on $X$.      
% \end{PROB}
% \begin{proof}
%     Since $\emptyset, X \in T_\alpha$ $\forall \alpha$, $\emptyset, X \in \bigcap \mathcal{T}_\alpha$. Moreover let $\left\{ \mathcal{B}_\beta \right\}$ be a generic finite subcollection of open sets such that $B_\beta \in \bigcap \mathcal{T}_\alpha$. By the definition of topology, $\bigcap B_\beta \in \mathcal{T}_\alpha$ for all $\alpha$. Thus $\bigcap B_\beta \in \bigcap \mathcal{T}_\alpha$. Similarly by the definition of topology, $\bigcup B_\beta \in \mathcal{T}_\alpha \implies \bigcup B_\beta \in \bigcap \mathcal{T}_\alpha$. Since $\left\{ B_\beta \right\}$ is a generic finite subcollection of $\bigcap \mathcal{T}_\alpha$, any statement is a valid statement for all finite subcollections. Thus $\bigcap \mathcal{T}_\alpha$ is a topology on $X$. 
    
%     Let $X = \R$. Let $T_1 := \left\{ \emptyset, \R, \R^+  \right\}$ and let $T_2 := \left\{ \emptyset, \R, \R^- \right\}$. $T_1 \cup T_2 = \left\{ 0, \R, \R^+, \R^- \right\}$. However $\R \setminus 0 = R^+ \cup R^- \not \in T_1 \cup T_2$. Thus $\bigcup T_\alpha$ is not a topology on $X$ by the existance of a contradiction      
% \end{proof}
% \begin{PROB}{4b}\label{Problem 4b.}
%     Let $\left\{ T_\alpha \right\}$ be a family of topologies on $X$. Show that there is unique smallest topology on X containing all $T_\alpha$, and a unique largest topology contained in all $T_\alpha$.     
% \end{PROB}
% \begin{proof}
%     Let $T$ be the topology generated by the subbasis $\bigcup T_\alpha$. For any topology $S$ which contains all $T_\alpha$. $\forall A \in T$, $A$ can be written as a finite intersection of elements of the subbasis $\bigcup T_\alpha$. Since $\bigcup T_\alpha \subseteq S \implies A \in S \implies T \subseteq S$. Thus $T$ is the unique smallest topology on X containing all $T_\alpha$. 
    
%     Let $T' = \bigcap \mathcal{T}_\alpha$. $T'$ is a topology (problem 4a) and is completely contained in all $T_\alpha$. Suppose $S$ is a topology of $X$ which is completely contained in all $T_\alpha$. Then $\forall A \in S, A \in \bigcap \mathcal{T}_\alpha \implies S \subseteq \bigcap \mathcal{T}_\alpha$. Thus $ \bigcap \mathcal{T}_\alpha$ is the largest topology contained by all $T_\alpha$.        
% \end{proof}
% \begin{PROB}{4c}\label{Problem 4c.}
%     If $X = \left\{ a, b, c \right\}$, let $T_1 = \left\{ \emptyset, X, a, ab \right\}$ and $T_2 = \left\{ \emptyset, X, a, bc \right\}$. Find the smallest topology containing $T_1$ and $T_2$, and the largest topology contained in $T_1$ and $T_2$.
% \end{PROB}
% \begin{proof}
%     $\left\{ \emptyset, X, a, ab, bc \right\}$ is the subbasis of the topology. $T = \left\{ \emptyset, X, a, b, ab, bc \right\}$ is the smallest topology containing $T_1$ and $T_2$. 
    
%     $S = \left\{ \emptyset, X, a \right\}$ is the largest topology contained by $T_1$ and $T_2$.   
% \end{proof}

% \bigskip

% \begin{PROB}{5}\label{Problem 5.}
%     Show that if A is a basis for a topology on X, then the topology generated by A equals the intersection of all topologies on X that contain A. Prove the same if A is a subbasis.
% \end{PROB}
% \begin{proof}
%     Let $T$ be the topology generated by basis A. Let $\left\{ S_\alpha \right\}$ be the collection of topologies on X which contain A. $\forall B \in T$, $B$ is a finite union of elements of the bassis A. Since $\forall \alpha$ $A \in S_\alpha$ and topologies are closed under union, $B \in S_\alpha$. Thus $B \in \bigcap S_\alpha \implies T \subseteq \bigcap S_\alpha$. $\forall B \in \bigcap S_\alpha$, thus $\forall \alpha$ $B \in S_\alpha$. Since $T = S_\alpha$ for some $\alpha$, $T$ is a topology which contains $A$. Thus $B \in T$. Thus $\bigcap S_\alpha \subseteq T$. Thus $T = \bigcap S_\alpha$.
    
%     Let $T$ be the topology generated by the subbasis A. Let $\left\{ S_\alpha \right\}$ be the collection of topologies on X which contain A. $T \in \left\{ S_\alpha \right\}$. $\forall B \in T$, $B$ is a finite union of finite intersections of elements of the subbasis A. Since $\forall \alpha$, $A \in S_\alpha$ and topologies are closed under intersection and unions. $B \in S_\alpha$. Thus $B \in \bigcap S_\alpha \implies T \subseteq \bigcap S_\alpha \implies T = \bigcap S_\alpha$.            
% \end{proof}

% \bigskip

% \begin{PROB}{6}\label{Problem 6.}
%     Show that $\R_l$ and $\R_K$ are not comparable. 
    
%     \[
%         \R_l := \left\{ [a, b) \mid a, b \in \R \right\}
%     \] 

%     \[
%         K := \left\{ \frac{1}{n} \mid n \in \N \right\}
%     \]

%     \[
%         \R_K := \left\{ (a, b) \mid a, b \in \R \right\} \cup \left\{ (a, b) \mid a, b \in \R \right\} \setminus K
%     \]
% \end{PROB}
% \begin{proof}
%     $[2, 3)$ is open in $\R_l$ but not open in $\R_k$. $\R \setminus K$ is open in $\R_k$ but not in $\R_l$. Thus $\R_l \not \subseteq \R_K$ and vice versa.        
% \end{proof}

% \bigskip

% \begin{PROB}{}\label{Problem .}
    
% \end{PROB}

\begin{claim}{2.1}\label{Claim 2.1.}
    Let $\left\{ U_i \right\}_{i = 1}^{n}$ be a finite collection of open sets in a topological space $(X, T)$. Then $\bigcap_{i = 1}^n U_i$ is open   
\end{claim}
\begin{proof}
    Base Case: $U_1 = \bigcap_{i = 1}^1 U_i$ is by definition a open set. 
    
    Inductive Hypothesis: Assume that for any $\left\{ U_i \right\}_{i = 1}^n$, finite collection of open sets, that $\bigcap_{i = 1}^n U_i$ is open. 
    
    Inductive Step: Let $\left\{ U_i \right\}_{i = 1}^{n + 1}$ is a finite collection of open sets. By the inductive hypothesis, $\bigcap_{i = 1}^n U_i$ is a open set. Thus $\bigcap_{i = 1}^{n + 1} U_i = \bigcap_{i = 1}^n \left(U_i\right) \cap U_{n + 1}$ is a open set. 
    
    
    Thus by induction, $\forall n \in \N$ if $\left\{ U_i \right\}_i^n$ is collection of open sets,  $\bigcap_{i = 1}^n U_i$ is a open set.  
\end{proof}

\bigskip

\begin{PROB}{2.2}\label{Problem 2.2.}
    Why does your proof not prove the false statement that the infinite intersection of open sets is necessarily open?
\end{PROB}
\begin{solution}
    Since my proof uses induction, it uses $n \in \N$ which necessitates that $n$ is finite. If $n$ was infinite the proof would be invalid as I am using induction to show that if the $n - 1$ case is true then that implies that the $n$ case is true.If $n$ was infinite, then the $n - 1$ and $n$ case would be the same and thus induction would not work.          
\end{solution}

\bigskip

\begin{claim}{2.3}\label{Claim 2.3.}
    A set U is open in a topological space $(X, T)$ if and only iff or every point $x \in U$ ,there exists an open set $U_x$ such that $x \in U_x \subset U$ .
\end{claim}
\begin{proof}
    $\implies$: Suppose that $U$ is a open set in $X$. Then $\forall x \in U$, let $U_x = U$. Thus $\exists U_x$ such that $U_x$ is open and $U_x \subseteq U$.
    
    $\impliedby$: Proven in problem 1 in Appendix.   
\end{proof}

\bigskip

\begin{PROB}{2.4}\label{Problem 2.4.}
    Verify that $\R_{s}^n$ (standard topology on $R^n$ open balls) is a topology on $\R^n$  
\end{PROB}
\begin{proof}
    (1) and (2) are trivially true. $\forall U, V \in \R_s^n$, by definition $\forall x \in U \cap V$, $\exists \delta_1, \delta_2 \ni B_{\delta_1}(x) \subseteq U, B_{\delta_2}(x) \subset V$. Let $\delta = \min(\delta_1, \delta_2)$. Thus $B_\delta(x) \subset U \cap V$. Thus $U \cap V$ is open by 2.2. Thus (3) is true. Let $\left\{ U_\alpha \right\}_{\alpha \in \lambda}$ be a finite collection of open sets. By definition, $\exists \delta_\alpha \ni B_{\delta_\alpha}(x) \subseteq U_\alpha \implies B_{\delta_\alpha}(x) \subseteq \bigcup_{\alpha \in \lambda}U_\alpha$. By 2.2, $\bigcup_{\alpha \in \lambda}U_\alpha$ is open. Thus $(4)$ is correct. Thus $\R_s^n$ is a valid topology over $\R^n$.         
\end{proof}

\bigskip

\begin{PROB}{2.6}\label{Problem 2.6.}
    Describe some of the open sets you get if $\R$  is endowed with the topologies described above (standard, discrete, indiscrete, finite complement, and countable complement). Specifically, identify sets that demonstrate the differences among these topologies, that is, find sets that are open in some topologies but not in others. For each of the topologies, determine if the interval $(0, 1)$ is an open set in that topology.
\end{PROB}
\begin{solution}
    All open sets in other topologies are in discrete. $\R \setminus \left\{ 0 \right\}$ is a open set in finite complement and countable complement but not a open set in standard and indiscrete. $\R \setminus \Q$ is a open set in countable complement but not in the other three. Finally any open interval is a openset in standard but not in the other three.
    
    \begin{enumerate}
        \item $(0, 1)$ is open in standard
        \item $(0, 1)$ is open in discrete
        \item $(0, 1)$ is not open in indiscrete
        \item $(0, 1)$ is not open in finite complement
        \item $(0, 1)$ is not open in countable complement     
    \end{enumerate}
    
\end{solution}

\bigskip

\begin{PROB}{2.7}\label{Problem 2.7.}
    Give an example of a topological space and a collection of opensets in that topological space that show that the infinite intersection of open sets need not be open.
\end{PROB}
\begin{solution}
    Let the topological space be $\R$ with a finite complement topology. $\left\{ \R \setminus \left\{ n \right\} \right\}_{n = 1}^{\infty}$ is a collection of open sets in this topological space. However $\bigcup_{n = 1}^{\infty} \R \setminus \left\{ n \right\} = \R \setminus \N$ is not open because $\N$ is not finite.     
\end{solution}

\section*{Limit Point and Closed Sets}
\begin{PROB}{2.8}\label{Problem 2.8.}
    Let $X = \R$ and let $A = (1, 2)$. Verify that 0 is a limit point in the indiscrete and finite complement topology but not in the standard topology or the discrete topology  
\end{PROB}
\begin{proof}
    $0$ is not in any open set in the indiscrete topology so it trivially satsifies the condition for a limit point. 

    Let $B \subset \R$ be a finite set where $0 \not \in B$. $\R \setminus B$ are the only open sets which 0 is in other than $\R$. $0 \in \R \setminus B$ and since $B$ is finite it can only remove a finite number of points in A, since A is uncountable, $(\R \setminus B) \cap A$ is not empty. Thus 0 is a limit point.
    
    $(-0.5, 0.5) \cap (1, 2) = \emptyset$, thus 0 is not a limit point in the standard topology. 
    
    $\left\{ 0 \right\}$ is a open set but its intersection with A is empty. So 0 is not a limit point in the discrete topology. 
\end{proof}

\bigskip

\begin{claim}{2.9}\label{Claim 2.9.}
    Suppose $p \not \in A$ in a topological space $(X, T)$. Then $p$ is not a limit point of $A$ if and only if there exists a neighborhood $U$ of $p$ such that $U \cap A = \emptyset$.  
\end{claim}
\begin{proof}
    $\implies:$ Suppose $p$ is not a limit point of $A$. Thus by the definition of a limit point, exists a open set $U \ni p \in U$ and $U \cap A = \emptyset$.
    
    $\impliedby:$ Suppose $\exists U \in T \ni p \in U, U \cap A = \emptyset$, by the definition of a limit point $p$ is not a limit point of $A$.     
\end{proof}

\bigskip

\begin{PROB}{2.10}\label{Problem 2.10.}
    If $p$ is an isolated point of a set A in a topological space X, show that there exists an open set U such that $U \cap A = \left\{ p \right\}$ .
\end{PROB}
\begin{proof}
    By definition since $p$ is a isolated point and not a limit point, $\exists U \in T \ni p \in U, (U \setminus {p}) \cap A = \emptyset$. But since $p \in A$ and $p \in U$, so $U \cap A = \left\{ p \right\}$.      
\end{proof}

\bigskip

% \begin{PROB}{2.11a}\label{Problem 2.11a.}
%     Examples of topological spaces and sets A where a limit point of A is an element of A. 
% \end{PROB}
% \begin{proof}
%     Let $X = \R$. 
    
%     \begin{enumerate}
%         \item Standard: Let $A = (0, 1)$ and let $p = 0.5$. 
%         \item Discrete Topology: No limit points
%         \item Indiscrete topology: $A = (0, 1)$, $\forall p \in (0, 1)$ is a limit point
%         \item Finite Complement Topology: $A = (0, 1)$. $\forall p \in (0, 1)$ is a limit point. Let $A = (0, 1) \in \Q$, only rationals. Then $\forall p \in (0, 1)$ is a limit point. 
%     \end{enumerate}
% \end{proof}
% \begin{PROB}{2.11b}\label{Problem 2.11b.}
%     Examples of topological spaces and sets A where a limit point of A is not an element of A. 
% \end{PROB}
% \begin{proof}
%     Let $X = \R$. 
    
%     \begin{enumerate}
%         \item Standard: Let $A = (0, 1)$ and let $p = 0$. 
%         \item Discrete Topology: No limit points
%         \item Indiscrete topology: $A = (0, 1)$, $\forall p \in \R \setminus A$ is a limit point
%         \item Finite Complement Topology: $A = (0, 1)$. $\forall p \in \R \setminus A$ is a limit point 
%     \end{enumerate}
% \end{proof}

\begin{PROB}{2.12a}\label{Problem 2.12a.}
    Which sets are closed in a set $X$ with the discrete topology? 
\end{PROB}
\begin{proof}
    All subsets of $X$ are closed with the discrete topology because there are no limit points in such a topology. 
\end{proof}
\begin{PROB}{2.12b}\label{Problem 2.12b.}
    Which sets are closed in a set $X$ with the indiscrete topology?
\end{PROB}
\begin{proof}
    Only $\emptyset$ and $X$ is closed in such a topology. Because the $\emptyset$ has no limit points, and the limit points for any $A \subset X$, are all of $X$. Thus $\overline{A} = X$. Thus the only closed sets are $\emptyset$ and $X$.        
\end{proof}
\begin{PROB}{2.12c}\label{Problem 2.12c.}
    Which sets are closed in a set X with the finite complement topology?
\end{PROB}
\begin{proof}
    If $X$ is finite then the problem reduces to the case of discrete topology. So suppose $X$ is not finite. The empty set, X, and any finite subset is closed under this topology,   
\end{proof}
\begin{PROB}{2.12d}\label{Problem 2.12d.}
    Which sets are closed in a set X with the countable complement topology?
\end{PROB}
\begin{proof}
    $\emptyset, X, $ any countable subset of $X$.   
\end{proof}

\bigskip

\begin{claim}{2.13}\label{Claim 2.13.}
    For any topological space $(X, T)$ and $A \subset X$, the set $\overline{A}$ is closed. That is, for any set A, $\overline{\overline{A}} = \overline{A}$.    
\end{claim}
\begin{proof}
    $\overline{A} \subseteq \overline{\overline{A}}$. $\forall x \in \overline{\overline{A}}$, then for all $U \subset T \ni x \in U$ such that $\exists y \in (U \setminus \left\{ x \right\}) \cap \overline{A}$. Thus $(U \setminus \left\{ y \right\}) \cap A \neq \emptyset \implies \left(U \setminus \left\{ x \right\}\right) \cap A \neq \emptyset \implies x \in \overline{A}$. Thus $\overline{\overline{A}} \subseteq \overline{A} \implies \overline{\overline{A}} = \overline{A}$. By definition of closed, $\overline{A}$ is closed.     
\end{proof}

\bigskip

\begin{claim}{2.14}\label{Claim 2.14.}
    Let $(X, T)$ be a topological space. Then the set $A$ is closed $\iff$ $X \setminus A$ is open.     
\end{claim}
\begin{proof}
    $\implies$: Suppose $A$ is closed. $\forall x \in X \setminus A$, since $\overline{A} = A \implies x$ is not a limit point of $A$. Thus $\exists$ a openset $U$ such that $x \in U, (U \setminus \left\{ x \right\}) \cap A = \emptyset \implies U \cap \left(X \setminus A\right) = U \implies U \subseteq X \setminus A$. Thus by Theorem 2.3, $X \setminus A$ is open. 
    
    $\impliedby$: Suppose $X \setminus A$ is open. Let $\forall x \in X \ni \forall U \in T \ni x \in U, \left(U \setminus \left\{ x \right\}\right) \cap A \neq \emptyset$, ie for all limit points of A. Assume the contrary, that $x \in X \setminus A \implies U \subset X \setminus A \text{ (Theorem 2.3) } \implies U \cap A= \emptyset$, which is a contradiction. Thus $x \in A$. Thus $\overline{A} = A \implies A$ is closed.       
\end{proof}

\bigskip

\begin{claim}{2.15}\label{Claim 2.15.}
    Let $(X, T)$ be a topological space, and let $U$ be a open set and $A$ be a closed set of X. Then the set $U - A$ is open and the set $A - U$ is closed      
\end{claim}
\begin{proof}
    $\forall x \in U - A$, $x \in U$ but not in $A$. Thus $\exists V_1 \in T \ni x \in V_1$, since $U$ is open, $V_1 \subseteq U$. Since $x \not \in A$ and since A is closed, $\exists V_2 \in T \ni x \in V_2 \ni (V_2 \setminus \left\{ x \right\})\cap A = \emptyset \implies V_2 \cap A = \emptyset$. Thus let $V = V_1 \cap V_2 \in T$, by closure of finite intersections. Morover $x \in V \implies V \neq \emptyset$ and $V \cap A = \emptyset$. But since $V \subset V_1 \implies V \subseteq U \implies V \subseteq U - A $. By Theorem 2.3, $U - A$ is open. 

    Alternate Proof: $U - A = U \cap A^c$. $A^c$ is open because A is closed. $U \cap A^c = U - A$ is open.   
    
    $A - U = A \cap U^c = (A^c \cap U)^c = (U \cap A^c)^c = (U - A)^c$. Since $U - A$ is open, $A - U$ is closed.     
\end{proof}

\bigskip

\begin{claim}{2.16}\label{Claim 2.16.}
    Let $(X, T)$ be a topological space. Then:
    
    \begin{enumerate}
        \item $\emptyset$ is closed
        \item $X$ is closed
        \item The union of finitely many closed sets is closed
        \item Let $\left\{ A_\alpha \right\}_{\alpha \in \lambda}$ be a collection of closed sets. Then $\bigcap_{\alpha \in \lambda} A_\alpha$ is closed    
    \end{enumerate}
\end{claim}
\begin{proof}
    \begin{enumerate}
        \item $X$ is open in X thus $X^c = \emptyset$ is closed by 2.14. 
        \item $\emptyset$ is open in $X$, thus $X \setminus \emptyset = X$ is closed by 2.14.
        \item $\forall n \in \N$, let $\left\{ A_k \right\}_{k = 1}^{n}$ be finite collection of closed sets. By Demorgan's Laws, $X \setminus \bigcup A_k = \bigcap (X \setminus A_k)$. Since $\forall k$ $A_k$ is closed, $X \setminus A_k$ is open by 2.14. 2.1, $X \setminus \bigcup A_k = \bigcap (X \setminus A_k)$ is open. By 2.14, $\bigcup A_k$ is closed. 
        \item Let $\left\{ A_\alpha \right\}_{\alpha \in \lambda}$ be a collection of closed sets. By 2.14, $\forall \alpha \in \lambda$ since $A_\alpha$ is closed $\implies$ $X \setminus A_\alpha$ is open. Thus $\bigcup_{\alpha \in \lambda} \left(X \setminus A_\alpha\right)$ is open. By 2.14 and Demorgan's Laws, $X \setminus \bigcup_{\alpha \in \lambda} \left(X \setminus A_\alpha\right) = \bigcap_{\alpha \in \lambda} \left(A_\alpha\right)$ is closed.               
    \end{enumerate}
\end{proof}

\bigskip

\begin{PROB}{2.17}\label{Problem 2.17.}
    Give an example to show that the union of infinitely many closed sets in a topological space may be a set that is not closed.
\end{PROB}
\begin{solution}
    Take the standard topology on $\R$. $\left\{ \left[\sum_{k = 1}^{n - 1} \frac{1}{k^2}, \sum_{k = 1}^{n} \frac{1}{k^2}\right] \right\}_{n = 1}^{\infty}$ is a infinite collection of closed sets. However $\frac{\pi^2}{6} = \lim_{n \to \infty} \sum_{k = 1}^{n} \frac{1}{k^2}$ is not in the union of the collection, because each interval is bounded above by $\sum_{k = 1}^{n} \frac{1}{k^2}$ which approaches $\frac{\pi^2}{6}$ but never equals it. Moreover $\forall \epsilon$, $\exists K \ni \forall n 
    \geq K$, $\abs{ \frac{\pi^2}{6} - \sum_{k = 1}^{n} \frac{1}{k^2} } < \epsilon$. Thus open sets containing $\frac{\pi^2}{6}$ have a non trivial intersection with the union of the the infinite collection. Thus $\frac{\pi^2}{6}$ is a limit point but is not in the union of the infinite collection.     
    
    
    Take finite complement topology on $\Z$. $\left\{ \left\{ n, n+1 \right\} \right\}_{n = 1}^{\infty}$ are sets of closed sets. The infinite union this is $\N$, but $\N$ is not closed in $\Z$ because it isn't finite.   
\end{solution}

\bigskip

\begin{PROB}{2.18a}\label{Problem 2.18a.}
    Topological Spaces and Sets which are are closed, but not open
\end{PROB}
\begin{proof}
    \begin{enumerate}
        \item $\R$ with countable complement topology. $\N, \Q, \R$ and any countable set. 
        \item $\Q$ with finite complement topology. $[0, \dots, 10]$ is closed but not open.       
    \end{enumerate}
\end{proof}
\begin{PROB}{2.18b}\label{Problem 2.18b.}
    Topological Spaces and Sets which are are open, but not closed
\end{PROB}
\begin{solution}
    \begin{enumerate}
        \item $\R$ and standard. $(a, b)$ is open but not closed
        \item $\R$ and discrete. Any subset of $\R$
        \item $\R$ and countable complement. Let $S$ be a countable subset of $\R$. $\R / S$ is open but not closed.        
    \end{enumerate}
\end{solution}
\begin{PROB}{2.18c}\label{Problem 2.18c.}
    Topological Spaces and Sets which are are open and closed
\end{PROB}
\begin{solution}
    \begin{enumerate}
        \item $\R$ and standard. Nothing is clopen. 
        \item $\Q$ and countable complement. $\Z$ is closed and open.         
    \end{enumerate}
\end{solution}
\begin{PROB}{2.18c}\label{Problem 2.18c.}
    Topological Spaces and Sets which are are neither open and closed
\end{PROB}
\begin{solution}
    \begin{enumerate}
        \item $\R$ and standard. $[a, b)$  
        \item $\R$ and countable complement. (0, 1) is neither open nor closed.        
    \end{enumerate}
\end{solution}

\bigskip

\begin{PROB}{2.19a}\label{Problem 2.19a.}
    Say if the sets are open closed or neither: 
    In $\Z$  with the finite complement topology: $\left\{ 0, 1, 2 \right\}$ , {prime numbers}, $\left\{ n \mid |n| \geq 10 \right\}$.
\end{PROB}
\begin{solution}
    \begin{enumerate}
        \item $A = \left\{ 0, 1, 2 \right\}$: $\Z \setminus A$ is infinite. Thus $A$ is not open. $\Z \setminus A$ is open because $A$ is finite. Thus $A$ is closed. 
        \item $A = \left\{ \text{prime numbers} \right\}$. $\Z \setminus A$ is infinite. Thus $A$ is not open. A is not finite thus $Z \setminus A$ is not open thus A is not closed. Thus A is neither open or closed.
        \item $A = \left\{ n \mid \abs{n} \geq 10 \right\}$. $Z \setminus A = \left\{ n \mid \abs{n} < 10 \right\}$ is finite. Thus A is open. A is infinite thus $Z \setminus A$ is not open thus A is not closed. Thus A is only open. 
    \end{enumerate}
\end{solution}
\begin{PROB}{2.19b}\label{Problem 2.19b.}
    Say if the sets are open, closed, both or neither: 
    In $\R$ with the standard topology: $(0, 1)$, $(0, 1]$, $[0, 1]$, $\left\{ 0, 1 \right\}$, $\left\{ \frac{1}{n} \mid n \in \N \right\}$  
\end{PROB}
\begin{solution}
    \begin{enumerate}
        \item $A = (0, 1)$: A is open by definition. 1 is a limit point because for all open sets $U$  which contain $1$, $\exists \delta \in \R^+ \ni (1 - \delta, 1 + \delta) \subseteq U$ and $(1 - \delta, 1 + \delta) \subseteq A \implies \left(U \setminus 1\right) \cap A \neq \emptyset$. Thus $1$ is a limit point but not in $A$. 
        \item $A = (0, 1]$. For all open sets U which $1$ is in, $\exists \delta \in \R^+ \ni (1 - \delta, 1 + \delta) \subseteq U$, but $(1 - \delta, 1 + \delta) \not \subset A \implies \forall U \in T \ni 1 \in U, U \not \subset A$. Thus A is not open. 0 is not in A. Thus A is neither open nor closed. 
        \item $A = [0, 1]$. Open sets which contain 0 and 1 are not subsets of A. Thus A is not open. $(-\infty, 0) \cup (1, \infty)$ is open thus A is closed. 
        \item $A = \left\{ 0, 1 \right\}$. Open sets which contain 0 or 1 are not subsets of A thus A is not open. A is closed. $(-\infty, 0) \cup (0, 1) \cup (1, \infty)$ is open thus A is closed. A has no limit points and thus $\overline{A} = A$. 
        \item $A = \left\{ \frac{1}{n} \mid n \in \N \right\}$. Not open. not closed.              
    \end{enumerate}
\end{solution}
\begin{PROB}{2.19c}\label{Problem 2.19c.}
    Say if the sets are open closed or neither: 
    In $\R^2$ with standard topology: 
    
    \begin{enumerate}
        \item $\left\{ (x, y) \mid x^2 + y^2 = 1 \right\}$ 
        \item $\left\{ (x, y) \mid x^2 + y^2 > 1 \right\}$
        \item $\left\{ (x, y) \mid x^2 + y^2 \geq 1 \right\}$
    \end{enumerate}
\end{PROB}
\begin{solution}
    \begin{enumerate}
        \item $A = \left\{ (x, y) \mid x^2 + y^2 = 1 \right\}$. Not open. Closed points on the ball
        \item $A = \left\{ (x, y) \mid x^2 + y^2 > 1 \right\}$. Open. Not closed
        \item $\left\{ (x, y) \mid x^2 + y^2 \geq 1 \right\}$. closed but not open. 
    \end{enumerate}
\end{solution}

\bigskip

\begin{claim}{2.20}\label{Claim 2.20.}
    For any set $A$ in a topological space $X$, the closure of $A$ equals the intersection of all closed sets containing $A$:
    
    \[
        \overline{A} = \bigcap_{A \subset B, B \in \mathcal{C}} B
    \]

    where $\mathcal{C}$ is the collection of closed sets in X. 
\end{claim}
\begin{proof}
    $D = \bigcap_{A \subset B, B \in \mathcal{C}} B \subseteq \overline{A}$ since $\overline{A}$ is a closed set and $A \subseteq \overline{A}$. 
    
    $\forall x \in \overline{A}$, $x \in A$ or $x$ is a limit point. If $x \in A$, then since $\forall B \in \mathcal{B} \ni A \subset B \implies x \in B \implies x \in D$. If $x$ is a limit point of $A$. Assume the contrary that $\exists B \in \mathcal{C} \ni A \subseteq B, x \not \in B$, since $B$ is closed then $X \setminus B$ is open. Since $x$ is a limit point of $A$ that means, $\left(X \setminus B\right) \cap A \neq \emptyset$. But since $\left(X \setminus B\right)\cap B = \emptyset$ and $A \subseteq B$, $\left(X \setminus B\right) \cap A = \emptyset$ which is a contradiction. Thus $x \in B$ for all $B$. Thus $x \in D$.               
    
    Thus $\overline{A} = D = \bigcap_{A \subset B, B \in \mathcal{C}} B$.  
\end{proof}

\bigskip

\begin{claim}{2.22a}\label{Claim 2.22a.}
    Let $A, B \subseteq X$ a topological space. Then $A \subseteq B \implies \overline{A} \subseteq \overline{B}$   
\end{claim}
\begin{proof}
    Let $A, B \subseteq X$ a topological space where $A \subseteq B$. $\forall x \in \overline{A}$, $x$ is a element of $A$ or a limit point of $A$. $A \subseteq B \subseteq \overline{B}$, hence if $x$ a element of $A$ then $x \in \overline{B}$. If $x$ is a limit point of $A$. Thus $\forall U \in T \ni x \in U, \left(U \setminus x\right)\cap A \neq \emptyset$. Since $A \subseteq B \implies \left(U \setminus x\right)\cap B \neq \emptyset$, thus $x$ is a limit point of $B$, thus $x \in \overline{B}$. Thus $\overline{A} \subseteq \overline{B}$. 
    
    \textbf{Alternate proof}: By definition $\overline{B} \in \mathcal{C}$. Moreover $A \subseteq B \subseteq \overline{B} \implies A \subseteq \overline{B}$. By claim 2.20, $\overline{A} \subseteq \overline{B}$ since $\overline{A} = \bigcap_{A \subseteq D, D \in \mathcal{C}} D$ and $\overline{B}$ is one such $D$.            
\end{proof}
\begin{claim}{2.22b}\label{Claim 2.22b.}
    Let $A, B \subseteq X$ a topological space. Then $\overline{A \cup B} = \overline{A} \cup \overline{B}$
\end{claim}
\begin{proof}
    $\forall x \in \overline{A \cup B}$, $x$ is a element of $A \cup B$ or $x$ is a limit point of $A \cup B$. If $x \in A \cup B \implies x \in A$ or $x \in B$. If $x \in A \implies x \in \overline{A} \implies x \in \overline{A} \cup \overline{B}$. Similarly with identical logic if $x \in B \implies x\in \overline{A} \cup \overline{B}$. If $x$ is a limit point of $A \cup B$, then $x$ is a limit point of $A$ or a limit point of $B$. If $x$ is a limit point of A then $x \in \overline{A} \implies x \in \overline{A} \cup \overline{B}$. If $x$ is a limit point of B then $x \in \overline{B}  \implies x \in \overline{A} \cup \overline{B}$. Thus $x \in \overline{A} \cup \overline{B} \implies \overline{A \cup B} \subseteq \overline{A} \cup \overline{B}$. If $x \in \overline{A} \cup \overline{B}$, $x \in \overline{A}$ or $x \in \overline{B}$. The logic for either case is identical so I will just prove that if $x \in \overline{A} \implies x \in \overline{A \cup B}$. $x \in \overline{A} \implies x\in A$ or $x$ is a limit point of $A$. If $x \in A \implies x \in A \cup B \implies x \in \overline{A \cup B}$. If $x$ is a limit point of $A$, then $\forall U \in T \ni x \in U, \left(U \setminus \left\{ x \right\}\right) \cap A \neq \emptyset \implies (U \setminus \left\{ x \right\}) \cap (A \cup B) \neq \emptyset$, $x$ is a limit point of $A \cup B$. Thus $x \in \overline{A \cup B}$. Thus $\overline{A} \cup \overline{B}\subseteq \overline{A \cup B} \implies \overline{A \cup B} = \overline{A} \cup \overline{B}$

\end{proof}

\bigskip

\begin{PROB}{2.23}\label{Problem 2.23.}
    Let $\left\{ A_{\alpha} \right\}_{\alpha \in \lambda}$ be a collection of subsets of a topological space $X$. Then is the statement
    
    \begin{equation*}
        \overline{\bigcup_{\alpha \in \lambda} A_{\alpha}} = \bigcup_{\alpha \in \lambda} \overline{A_\alpha}
    \end{equation*}
    true?
\end{PROB}
\begin{solution}
    If $\left\{ A_{\alpha} \right\}_{\alpha \in \lambda}$ is a finite collection this is automatically true by Theorem 2.22b. Now for infinite collections it is saying the infinite union of closed sets is closed which we proved a counter example for in 2.17. 
\end{solution}

\bigskip

\begin{PROB}{2.24}\label{Problem 2.24.}
    In $\R^2$ with the standard topology, describe the limit points and closure of each of the following two sets:
    
    \begin{enumerate}
        \item Topologist's sine curve: $S = \left\{ (x, \sin\left(\frac{1}{x}\right)) \mid x \in (0, 1) \right\}$
        \item Topologist's comb: $C = \left\{ (x, 0) \mid x \in [0, 1] \right\} \cup \bigcup_{n = 1}^{\infty} \left\{ (\frac{1}{n}, y) \mid y \in [0, 1] \right\}$  
    \end{enumerate}
    
\end{PROB}
\begin{solution}
    \begin{enumerate}
        \item The limit points of $S$ are $S$, $\left\{ (0, y) \mid y \in [-1, 1] \right\}$, and $\left\{ (1, \sin(1) \right\}$ . Thus $\overline{S} = S \cup \left\{ (0, y) \mid y \in [-1, 1] \right\} \cup \left\{ (1, \sin(1)) \right\}$. 
        \item The limit points of $T$ are $T$ and $\left\{ 0 \right\} \times [0, 1]$. Thus $\overline{T} = T \cup (\left\{ 0 \right\} \times [0, 1])$      
    \end{enumerate}
\end{solution}

\bigskip

\begin{PROB}{2.25}\label{Problem 2.25.}
    In the standard topology of $\R$, there exists a nonempty subset $C$ of $[0, 1]$ that is closed, contains no nonempty open interval, and no isolated points.   
\end{PROB}
\begin{proof}
    Let $A_0 = [0, 1]$. Let $A_{i + 1}$ be the union of the first third and last third of each closed interval of $A_i$. Example: $A_1 = [0, \frac{1}{3}] \cup [\frac{2}{3}, 1]$, $A_2 = [0, \frac{1}{9}] \cup [\frac{2}{9}, \frac{1}{3}] \cup [\frac{2}{3}, \frac{7}{9}] \cup [\frac{8}{9}, 1]$. Then let $C = \bigcap_{n = 0}^{\infty} A_n$. $C$ is closed. $measure(C) \to 0$, no nonempty open intervals. But each point in $C$ is a limit point.        
\end{proof}

\pagebreak

\section*{Interior Points and Boundary Points}
\begin{claim}{2.26}\label{Claim 2.26.}
    Let $A \subseteq X$. Then $p \in A^{\circ}$ (Interior of $A$) iff $\exists U$, a open set, such that $p \in U \subseteq A$.    
\end{claim}
\begin{proof}
    $\implies$: Suppose $p$ is a interior point. Then by definition $p \in A^\circ$, which is a open set such that $A^\circ \subseteq A$.
    
    $\impliedby:$ Suppose $\exists U \in T \ni p \in U \subseteq A$, since $U \subseteq A \implies U \subseteq A^\circ \implies p \in A^\circ \implies p$ is a interior point.   
\end{proof}

\bigskip

\begin{PROB}{2.27}\label{Problem 2.27.}
    Show that a set $U$ is open in a topological space $X$ iff $\forall p \in U$, $p \in U^\circ$, in otherwords if p is a interior point.     
\end{PROB}
\begin{proof}
    $\implies$: U is open. By 2.3, $\forall p \in U, \exists U_p \in T \ni p \in U_p \subseteq U$. Since $U_p \subseteq U, U_p \subseteq U^{\circ}$, by definition of union. Thus $p \in U^\circ$. 
    
    $\impliedby:$ Suppose $\forall p \in U, p$ is a interior point. Thus $p \in U^\circ$ a open set where $U^\circ \subseteq U$. Thus by 2.3, U is open.    
\end{proof}

\bigskip

\begin{claim}{2.28a}\label{Claim 2.2a.}
    Let $A \subseteq X$, a topological space. Then $A^\circ, \delta A, (X - A)^\circ$ are disjoint.    
\end{claim}
\begin{proof}
    $\forall x \in A^\circ$, then $\exists U_x \in T \ni x \in U_x \subseteq A^\circ \subseteq A \implies U_x \subseteq A \implies U_x \cap \left(X \setminus A\right) = \emptyset$. Thus $U_x \cap (X - A)^\circ = \emptyset$. Assume the contrary, $x \in (X - A)^\circ$. Thus $\exists V_x \in T \ni x \in V_x \subseteq (X - A)^\circ$. Thus $\emptyset \neq U_x \cap V_x \subseteq A$ and $U_x \cap V_x \subseteq X - A$, which is a contradiction. Thus $x \not \in (X - A)^\circ$. Assume the contrary that $x \in \delta A$, thus $x \in \overline{A} \cap \overline{X - A} \implies x \in \overline{X - A}$. If $x \in X - A$, it produces a contradiction since $x \in A$ and $x \in X -A$. Thus if $x \in \overline{X -A}$, it must be a limit point, however $U_x \subseteq A \implies U_x \cap \left(X \setminus A\right) = \emptyset$, which is a contradiction. Thus $x \not \in \delta A$. Thus $A^\circ \cap (X - A)^\circ = \emptyset$ and $A^\circ \cap \delta A = \emptyset$.

    $\forall x \in \delta A = \overline{A} \cap \overline{X - A}\implies x \in \overline{A}$. Thus $x \in A$ or $x$ is a limit point of $A$. If $x \in A$ then $x \not \in X - A \implies x \not \in (X - A)^\circ$. Assume the contrary that $x \in (X - A)^{\circ}$ and $x$ is a limit point of $A$, then $\exists U \in T \ni x \in U \subset (X - A)^\circ \subseteq X - A \implies U \cap A = \emptyset$ which contradicts the fact that $x$ is a limit point of A and $(U \setminus \left\{ x \right\}) \cap A \neq \emptyset$. Thus by contradiction if $x$ is a limit point $x \not \in (X - A)^\circ$. Thus $x \not \in (X - A)^\circ$. Thus $\delta A \cap (X - A)^\circ = \emptyset$. 
    
    Thus $A^\circ, \delta A, (X - A)^\circ$ are disjoint. 
\end{proof}
\begin{claim}{2.28b}\label{Claim 2.28b.}
    Let $A \subseteq X$, a topological space. Then $A^\circ \cup \delta A \cup (X - A)^\circ = X$. 
\end{claim}
\begin{proof}
    $A^\circ \cup \delta A \cup (X - A)^\circ \subseteq X$. $\forall x \in X$, $x \in A$ or $x \in X - A$. 
    
    \begin{enumerate}
        \item If $x \in A$: If $\exists U \in T \ni x \in U, (U \setminus \left\{ x \right\}) \cap (X - A) = \emptyset \implies U \subseteq A \implies x \in U \subseteq A^\circ$. Otherwise $\forall U \in T \ni x \in U \implies (U \setminus \left\{ x \right\}) \cap (X - A) \neq \emptyset \implies x \in \overline{X - A}$. Since $x \in A \implies x \in \overline{A}$. Thus $x \in \overline{A} \cap \overline{X - A} = \delta A$. Thus if $x \in A \implies x \in A^\circ$ or $x \in \delta A$.         
        \item If $x \in X - A$: If $\exists U \in T \ni x \in U, (U \setminus \left\{ x \right\}) \cap A = \emptyset \implies U \subseteq X - A \implies x \in U \subseteq (X - A)^\circ$. Otherwise that means $\forall U \in T \ni x \in U \implies \left(U \setminus {x
        }\right) \cap A \neq \emptyset \implies x \in \overline{A}$. Since $x \in X - A \subseteq \overline{X - A}$. Thus $x \in \delta A$. Thus if $x \in X - A \implies x \in (X - A)^\circ$ or $x \in \delta A$. 
    \end{enumerate}

    Thus $x \in A^\circ$, $x \in \delta A$, or $x \in \overline{X - A} \implies x \in A^\circ \cup \delta A \cup (X - A)^\circ$. Thus $X \subseteq A^\circ \cup \delta A \cup (X - A)^\circ \implies A^\circ \cup \delta A \cup (X - A)^\circ = X$.    
    
\end{proof}

\bigskip

\begin{PROB}{2.29}\label{Problem 2.29.}
    Pick several different subsets of $\R$, and for each one, find the interior and the boundary using:

    \begin{enumerate}
        \item Discrete topology
        \item Indiscrete topology
        \item Finite Complement topology
        \item Standard Topology
    \end{enumerate}
\end{PROB}
\begin{proof}
    Analyze the following subsets of $\R$: $(0, 1]$, $C$ in problem 2.25, $\Q$, and $\left\{ 0, \dots, 10 \right\}$. 


    \begin{enumerate}
        \item Standard Topology
        \begin{enumerate}
            \item $A = (0, 1]$. $A^\circ = (0, 1)$. $\delta A = \left\{ 0, 1 \right\}$ 
            \item $C$. $C^\circ = \emptyset$. $\delta C = C$. 
            \item $\Q$. $\Q^\circ = \emptyset$. $\delta \Q = \R$
            \item $A = \left\{ 0, \dots, 10 \right\}$. $A^\circ = \emptyset$, $\delta A = A$           
        \end{enumerate}
        \item Discrete Topology

        \item Indiscrete topology
        \begin{enumerate}
            \item $A = (0, 1]$. $A^\circ = \emptyset$. $\delta A = \R$
                
        \end{enumerate}
    \end{enumerate}
    
    
\end{proof}

\pagebreak

\section*{Convergence of Sequences}

\begin{claim}{2.30}\label{Claim 2.30.}
    Let $A$ be a subset of a topological space $X$, and let $p \in X$. If the set $\left\{ x_i \right\}_{i \in \N} \subseteq A$ and $x_i \to p \implies p \in \overline{A}$     
\end{claim}
\begin{proof}
    By definition of limit of sequence, since $x_i \to p \implies \forall U \in T \ni p \in U, \exists K \in \N \ni \forall n \geq K, x_n \in U \implies x_i \in \left(U \setminus x\right) \cap A$. Thus $x_i$ is a limit point of $A$ thus $x \in \overline{A}$    
\end{proof}

\bigskip

\begin{claim}{2.31}\label{Claim 2.31.}
    In the standard topology on $\R^n$, if $p$ is a limit point of a set $A$, then there is a sequence $\left\{ x_n \right\}_{n \in \N} \subseteq A$ such that $x_n \to p$.     
\end{claim}
\begin{proof}
    Suppose $p$ is a limit point of $A$. Let us create this sequence by construction. First create a sequence of open sets. Let $B_n = B(p, \frac{1}{n})$ be a ball around $p$ of radius $\frac{1}{n}$. Note $B_1 \supset B_2 \supset \dots \supset B_i \supset \dots$ and $\frac{1}{n} \to 0$. Since p is a limit point of $A$, $\forall i$ $\left(B_i \setminus \left\{ p \right\}\right) \cap A \neq \emptyset$. Pick $x_i \in \left(B_i \setminus \left\{ p \right\}\right) \cap A$. By definition of standard topology, $\forall U \in T \ni p \in U$, $\exists \epsilon \ni B(p, \epsilon) \subseteq U$. By the archemedian property of reals, $\exists K \in \N \ni K > \frac{1}{\epsilon} \implies \epsilon > \frac{1}{K}$. Thus $\forall n \geq K, B_n \subseteq B(p, \epsilon) \subseteq U \implies x_n \in U$. Thus $p$ is the limit of the sequence $x_i$.               
\end{proof}

\bigskip

\begin{PROB}{2.32}\label{Problem 2.32.}
    Find an example of a topological space and a convergent sequence in that space for which the limit of the sequence is not unique.
\end{PROB}
\begin{proof}
    Let the $X$ be any set and let the topology be indiscrete topology. Let $x_i$ be any sequence. Then $\forall p \in X$, p is a limit of $x_i$. 
\end{proof}

\bigskip


\pagebreak

\section*{Appendix}
\begin{PROB}{1}\label{Problem 1.}
    Let $X$ be a topological space. Let $A \subseteq X$. $\forall x \in A$, there is a open set $U$ containing $x$ such that $U \subset A$. Show that $A$ open in $X$        
\end{PROB}
\begin{proof}
    $\forall x \in A$, define $U_x$ as the open set containing $x$ such that $U_x \subset A$. Let $B := \bigcup_{x \in A} U_x$. 
    
    $\forall x \in A, x \in U_x$ by definition. $U_x \subseteq \bigcup_{x \in A} U_x$. Thus $x \in B$. Thus $A \subseteq B$.
    
    $\forall x \in B$, by definition of union, $\exists U_x$, open set such that $x \in U_x$. Since $U_x \subset A$, thus $x \in A$. Thus $A \subseteq B$.
    
    Thus $A = B$. 
\end{proof}

\bigskip

\end{document}